\begin{table}[H]
\centering
\caption{Baseline phenotyping method comparison in the 2022--2023 temporal-validation cycle. Mean maximum prevalence ratio is averaged across prespecified endpoints. Minimum cluster proportion and bootstrap ARI are development/test stability summaries; a high prevalence ratio from a very small cluster should not be interpreted as a stronger method.}
\label{tab:baseline_methods}
\scriptsize
\setlength{\tabcolsep}{3pt}
\begin{tabular*}{\textwidth}{@{\extracolsep{\fill}}>{\raggedright\arraybackslash}p{0.32\textwidth}rrrr@{}}
\toprule
Method & Mean max PR & Boot. ARI & Min. cluster, \% & Endpoints \\
\midrule
MCA-style one-hot SVD + k-means & 1.57 & 0.774 & 27.5 & 5 \\
Raw matrix PCA + k-means & 2.25 & 0.735 & 0.1 & 5 \\
SSL embedding + k-means & 2.16 & 0.953 & 4.9 & 5 \\
Selected-variable Bernoulli LCA & 1.50 & 0.904 & 18.5 & 5 \\
\bottomrule
\end{tabular*}
\end{table}
